% !TeX root = ../../book.tex
\begin{chapex}
判断下列方程是否存在函数 $f : \mathbb{R} \to \mathbb{R}$ 使得，对于所有 $x,y \in \mathbb{R}$，方程成立当且仅当 $y=f(x)$。
\begin{multicols}{3}
\begin{enumerate}[(a)]
\item $x+y=1$
\item $x^2+y^2=1$
\item $x=0$
\item $y=0$
\item $(1+x^2)y=1$
\item $(1-x^2)y=0$
\end{enumerate}
\end{multicols}
\end{chapex}

\begin{chapex}
设 $X$ 为集合。证明。
\[ \forall a \in X,\, \exists ! U \in \mathcal{P}(X), (a \in U \wedge \exists ! x \in X,\, x \in U) \]
给出该逻辑公式所给出函数 $X \to \mathcal{P}(X)$ 的明确描述。
\hintlabel{cqSingletonSubsetFunction}{%
}
\end{chapex}

\begin{chapex}
证明只有一个函数的值域为空。它的定义域是什么？
\hintlabel{cqFunctionEmptyCodomain}{%
给定一个函数 $f : X \to \varnothing$，对于每个 $a \in X$，我们必有 $f(a) \in \varnothing$。
}
\end{chapex}

\begin{definition}
\label{defEvenOddFunction}
\index{even!function}
\index{odd!function}
函数 $f : \mathbb{R} \to \mathbb{R}$ 如果对于所有 $x \in \mathbb{R}$，有 $f(-x)=f(x)$，则称 $f$ 为\textbf{偶函数}，如果对于所有 $x \in \mathbb{R}$，有 $f(-x)=-f(x)$，则称 $f$ 为\textbf{奇函数}。
\end{definition}

\begin{chapex}
设 $n \in \mathbb{N}$。函数 $f : \mathbb{R} \to \mathbb{R}$ 定义为对所有 $x \in \mathbb{R}$, $f(x)=x^n$。证明函数 $f$ 为偶函数当且仅当 $n$ 为偶数，$f$ 为奇函数当且仅当 $n$ 为奇数。
\end{chapex}

\begin{chapex}
证明只有唯一一个函数 $f : \mathbb{R} \to \mathbb{R}$ 既是奇函数又是偶函数。
\end{chapex}

\begin{chapex}
设 $U \subseteq \mathbb{R}$，并设 $\chi_U : \mathbb{R} \to \{0,1\}$ 为 $U$ 的指示函数。
\begin{enumerate}[(a)]
\item 证明 $\chi_U$ 为偶函数当且仅当 $U = \{ -u \mid u \in U \}$；
\item 证明 $\chi_U$ 为奇函数当且仅当 $\mathbb{R} \setminus U = \{ -u \mid u \in U \}$。
\end{enumerate}
\end{chapex}

\begin{chapex}
证明对于所有函数 $f : \mathbb{R} \to \mathbb{R}$，存在唯一偶函数 $g : \mathbb{R} \to \mathbb{R}$ 和唯一奇函数 $h : \mathbb{R} \to \mathbb{R}$ 使得对于所有 $x \in \mathbb{R}$，都有 $f(x)=g(x)+h(x)$。
\hintlabel{cqEveryFunctionIsUniquelySumOfEvenAndOdd}{%
考虑对于 $x \in \mathbb{R}$, $f(x)+f(-x)$ 和 $f(x)-f(-x)$。
}
\end{chapex}

\begin{chapex}
设 $\{ \theta_n : [n] \to [n] \mid n \in \mathbb{N} \}$ 为函数族，满足对于所有 $f : [m] \to [n]$， $f \circ \theta_m = \theta_n \circ f$。证明对于所有 $n \in \mathbb{N}$, $\theta_n = \mathrm{id}_{[n]}$。
\hintlabel{cqNaturalTransformationFromIdToId}{%
给定 $n \in \mathbb{N}$ 并思考 $\theta_n$ 如何与函数 $f : [1] \to [n]$ 交互。
}
\end{chapex}

\begin{chapex}
设 $X$ 为集合并设 $U, V \subseteq X$。用 $\chi_U$ 和 $\chi_V$ 描述 $U$ 和 $V$ 对称差的指示函数 $\chi_{U \symmdiff V}$（参见\Cref{defSymmetricDifference}）。
\end{chapex}

\begin{chapex}
\label{cqAntiderivativeOfLogarithm}
[本问题假设你对积分有一定了解。] 微积分学生经常听到的一个谎言是
\[ \int \dfrac{1}{x} \, dx = \log(|x|) + c \]
其中 $\log$ 表示自然对数函数，$c$ 是任意实数常数。证明正确积分公式如下
\[ \int \dfrac{1}{x} \, dx = \log(|x|) + a \chi_U(x) + b \chi_V(x) \]
其中 $a$ 和 $b$ 为任意实数常数，且 $U$ 和 $V$ 为由你确定的 $\mathbb{R} \setminus \{ 0 \}$ 的特定（非空）子集。
\end{chapex}

\subsection*{像与原像}

\Crefrange{cqComputeImageBegin}{cqComputeImageEnd}，找出问题描述函数 $f$ 定义域子集 $U$ 的像 $f[U]$。

\begin{chapex}
\label{cqComputeImageBegin}
$f : \mathbb{R} \to \mathbb{R}$；对于所有 $x \in \mathbb{R}$; $U = \mathbb{R}$, $f(x) = \sqrt{1+x^2}$ 。
\end{chapex}

\begin{chapex}
$f : \mathbb{Z} \times \mathbb{Z} \to \mathbb{Z}$；对于所有 $(a,b) \in \mathbb{Z} \times \mathbb{Z}$, $f(a,b) = a+2b$; $U = \{ 1 \} \times \mathbb{Z}$。
\end{chapex}

\begin{chapex}
$f : \mathbb{N} \to \mathcal{P}(\mathbb{N})$；对于所有 $n \in \mathbb{N}$, $f(0) = \varnothing$ 且 $f(n+1) = f(n) \cup \{ n \}$; $U = \mathbb{N}$。
\end{chapex}

\begin{chapex}
$f : \mathbb{R}^{\mathbb{R}} \to \mathbb{R}^{\mathbb{R}}$ （其中 $\mathbb{R}^{\mathbb{R}}$ 为所有 $\mathbb{R} \to \mathbb{R}$ 函数的集合）；对于所有 $h \in \mathbb{R}^{\mathbb{R}}$ 和所有 $x \in \mathbb{R}$, $f(h)(x) = h(|x|)$; $U = \mathbb{R}^{\mathbb{R}}$.
\hintlabel{cqComputeImageEnd}{%
首先观察每个 $h \in f[\mathbb{R}^{\mathbb{R}}]$ 在 \Cref{defEvenOddFunction} 意义上为偶函数。
}
\end{chapex}

\Crefrange{cqComputePreimageBegin}{cqComputePreimageEnd}，找出问题描述函数 $f$ 值域子集 $V$ 的原像 $f^{-1}[V]$。

\begin{chapex}
\label{cqComputePreimageBegin}
$f : \mathbb{R} \to \mathbb{R}$；对于所有 $x \in \mathbb{R}$, $f(x) = \sqrt{1+x^2}$; $V = (-5,5]$。
\end{chapex}

\begin{chapex}
$f : \mathbb{Z} \times \mathbb{Z} \to \mathbb{Z}$；对于所有 $(a,b) \in \mathbb{Z} \times \mathbb{Z}$, $f(a,b) = a+2b$; $V = \{ n \in \mathbb{Z} \mid n \text{为奇数} \}$。
\end{chapex}

\begin{chapex}
\label{cqComputePreimageEnd}
$f : \mathcal{P}(\mathbb{N}) \times \mathcal{P}(\mathbb{N}) \to \mathcal{P}(\mathbb{N})$；对于所有 $(A,B) \in \mathcal{P}(\mathbb{N}) \times \mathcal{P}(\mathbb{N})$, $f(A,B) = A \cap B$; $V = \{ \varnothing \}$。
\end{chapex}

\begin{chapex}
设 $f : X \to Y$ 为函数。对于以下陈述，证明其为真或给出一个反例。
\begin{multicols}{2}
\begin{enumerate}[(a)]
\item 对于所有 $U \subseteq X$, $U \subseteq f^{-1}[f[U]]$;
\item 对于所有 $U \subseteq X$, $f^{-1}[f[U]] \subseteq U$;
\item 对于所有 $V \subseteq Y$, $V \subseteq f[f^{-1}[V]]$;
\item 对于所有 $V \subseteq Y$, $f[f^{-1}[V]] \subseteq V$。
\end{enumerate}
\end{multicols}
\end{chapex}

\begin{chapex}
设 $f : X \to Y$ 为函数，设 $A$ 为集合，并设 $p : X \to A$ 且 $i : A \to Y$ 为函数使得如下条件成立：
\begin{enumerate}[(i)]
\item $i$ 为单射；
\item $i \circ p = f$；
\item 如果 $q : X \to B$ 和 $j : B \to Y$ 为函数，满足 $j$ 为单射且 $j \circ q = f$，则存在唯一函数 $u : A \to B$ 使得 $j \circ u = i$。
\end{enumerate}
证明存在唯一双射 $v : A \to f[X]$ 使得对于所有 $a \in f[X]$, $i(a)=v(a)$。
\hintlabel{cqUniversalPropertyOfImage}{%
这个问题非常棘手。首先证明当 $A = f[X]$ 且适当选择 $p$ 和 $i$ 时，满足条件 (i)--(iii)。然后条件 (iii) 在每种情况下（对于 $A$ 和 $f[X]$）定义函数 $v : A \to f[X]$ 和 $w : f[X] \to A$，并给出 $v$ 的唯一性。您可以使用条件 (iii) 的`唯一性'部分证明这些函数是互逆的。
}
\end{chapex}

\begin{chapex}
设 $f : X \to Y$ 为函数并设 $U, V \subseteq Y$。证明：
\begin{enumerate}[(a)]
\item $f^{-1}[U \cap V] = f^{-1}[U] \cap f^{-1}[V]$;
\item $f^{-1}[U \cup V] = f^{-1}[U] \cup f^{-1}[V]$; 
\item $f^{-1}[Y \setminus U] = X \setminus f^{-1}[U]$。
\end{enumerate}
因此原像保留了基本的集合运算。
\end{chapex}

\begin{chapex}
设 $f : X \to Y$ 和 $g : Y \to Z$ 为函数。
\begin{enumerate}[(a)]
\item 证明对于所有 $U \subseteq X$, $(g \circ f)[U] = g[f[U]]$;
\item 证明对于所有 $W \subseteq Z$, $(g \circ f)^{-1}[W] = f^{-1}[g^{-1}[W]]$。
\end{enumerate}
\end{chapex}

\subsection*{单射、满射和双射}

\begin{chapex}
\begin{enumerate}[(a)]
\item 证明对于所有函数 $f : X \to Y$ 和 $g : Y \to Z$，如果 $g \circ f$ 为双射，则 $f$ 为单射且 $g$ 为满射。
\item 找到函数 $f : X \to Y$ 和函数 $g : Y \to Z$ 的一个例子，满足 $g \circ f$ 为双射，$f$ 不是满射且 $g$ 不是单射。
\end{enumerate}
\hintlabel{cqCompositeOfFunctionsIsBijection}{%
避免通过矛盾来证明此问题的任何一部分。对于（a），可以直接通过`单射'和`满射'的定义获得简短的证明。对于 (b)，找到尽可能简单的反例。
}
\end{chapex}

\begin{chapex}
\label{cqInjectionSurjectionBijection}
对于下面 $\mathbb{R}$ 的每个子集对 $(U,V)$，判断 `对于所有 $x \in U$, $f(x) = x^2-4x+7$' 是否定义函数 $f : U \to V$，如果是，判断 $f$ 是单射还是满射。
\begin{multicols}{2}
\begin{enumerate}[(a)]
\item $U = \mathbb{R}$ 且 $V = \mathbb{R}$;
\item $U = (1, 4)$ 且 $V = [3, 7)$;
\item $U = [3, 4)$ 且 $V = [4, 7)$;
\item $U = (3, 4]$ 且 $V = [4, 7)$;
\item $U = [2, \infty)$ 且 $V = [3, \infty)$;
\item $U = [2,\infty)$ 且 $V = \mathbb{R}$.
\end{enumerate}
\end{multicols}
\end{chapex}

\begin{chapex}
对于下面每组集合 $X$ 和 $Y$，找到（并证明）双射 $f : X \to Y$。
\begin{enumerate}[(a)]
\item $X = \mathbb{Z}$ 且 $Y = \mathbb{N}$;
\item $X = \mathbb{R}$ 且 $Y = (-1,1)$;
\item $X = [0,1]$ 且 $Y = (0,1)$;
\item $X = [a,b]$ 且 $Y = (c,d)$, 其中 $a,b,c,d \in \mathbb{R}$ 且 $a<b$, $c<d$。
\end{enumerate}
\end{chapex}

\begin{chapex}
证明函数 $f : \mathbb{N} \times \mathbb{N} \to \mathbb{N}$ 定义为对于所有 $(a,b) \in \mathbb{N} \times \mathbb{N}$, $f(a,b) = \dbinom{a+b+1}{2} + b$ 为双射。
\hintlabel{cqElementaryBijectionNTimesNToN}{%
首先证明对于所有 $m \ge 1$, $\dbinom{m}{2} < \dbinom{m+1}{2}$。推出对于所有 $n \in \mathbb{N}$，存在唯一自然数 $k$ 满足 $\dbinom{k+1}{2} \le n < \dbinom{k+2}{2}$。你能看出这与函数 $f$ 有什么关系吗？
}
\end{chapex}

\begin{chapex}
设 $e : X \to X$ 为函数且满足 $e \circ e = e$。证明存在集合 $Y$ 和函数 $f : X \to Y$ 与 $g : Y \to X$ 满足 $g \circ f = e$ 且 $f \circ g = \mathrm{id}_Y$.
\hintlabel{exSplitIdempotents}{%
考虑 $e$ 的不动点集合 --- 即元素 $x \in X$ 使得 $e(x)=x$。
}
\end{chapex}

\subsection*{判断题}

\tfquestiontext{cqFunctionsTFBegin}{cqFunctionsTFEnd}

\begin{chapex} % False
\label{cqFunctionsTFBegin}
集合 $G = \{ (x,y) \in \mathbb{R} \times \mathbb{R} \mid x^2=y^2 \}$ 是某函数的图像。
\end{chapex}

\begin{chapex} % True
集合 $G = \{ (x,y) \in \mathbb{Z}\times \mathbb{N} \mid x^2=y^2 \}$ 是某函数的图像。
\end{chapex}

\begin{chapex} % False
每个具有空定义域的函数都有一个空值域。
\end{chapex}

\begin{chapex} % True
每个具有空值域的函数都有一个空定义域。
\end{chapex}

\begin{chapex} % False (in general)
函数的像是其定义域的子集
\end{chapex}

\begin{chapex} % True
给定函数 $f : X \to Y$，赋值 $V \mapsto f^{-1}[V]$ 定义函数 $\mathcal{P}(Y) \to \mathcal{P}(X)$。
\end{chapex}

\begin{chapex} % True
设 $f$、$g$ 和 $h$ 为可组合函数。如果 $h \circ g \circ f$ 为单射，则 $g$ 为单射。
\end{chapex}

\begin{chapex} % True
\label{cqFunctionsTFEnd}
所有左逆都是满射，所有右逆都是单射。
\end{chapex}

\subsection*{可能性问题}

\asnquestiontext{cqFunctionsASNBegin}{cqFunctionsASNEnd}

\begin{chapex} % Always
\label{cqFunctionsASNBegin}
设 $f : X \to Y$ 为函数，并设 $U \subseteq X$。则存在函数 $g : U \to Y$ 定义为对于所有 $x \in U$, $g(x) = f(x)$。
\end{chapex}

\begin{chapex} % Sometimes
设 $f : X \to Y$ 为函数，并设 $V \subseteq Y$。则存在函数 $g : X \to V$ 定义为对于所有 $x \in X$, $g(x) = f(x)$。
\end{chapex}

\begin{chapex} % Always
设 $X$ 为集合。则存在唯一函数 $X \to \{ 0 \}$。
\end{chapex}

\begin{chapex} % Sometimes
设 $X$ 为集合。则存在唯一函数 $X \to \varnothing$。
\end{chapex}

\begin{chapex} % Sometimes
设 $X$ 和 $Y$ 为集合，并设 $G \subseteq X \times Y$。则 $G$ 为函数 $f : X \to Y$ 的图像。
\end{chapex}

\begin{chapex} % Never
设 $f : \{ 1, 2, 3 \} \to \{ 1, 2 \}$，设 $G = \mathrm{Gr}(f)$，设 $G' = \{ (y,x) \mid (x,y) \in G \}$。则 $G'$ 是函数 $\{ 1, 2 \} \to \{ 1, 2, 3 \}$ 的图像。
\end{chapex}

\begin{chapex} % Always
设 $f : X \to Y$ 为函数，并设 $U \subseteq X$ 非空。则 $f[U]$ 非空。
\end{chapex}

\begin{chapex} % Sometimes
设 $f : X \to Y$ 为函数，并设 $V \subseteq Y$ 非空。则 $f^{-1}[V]$ 非空。
\end{chapex}

\begin{chapex} % Sometimes
设 $f : \{ 1, 2 \} \to \{ 1, 2, 3 \}$ 为函数。则 $f$ 为单射。
\end{chapex}

\begin{chapex} % Never
设 $f : \{ 1, 2 \} \to \{ 1, 2, 3 \}$ 为函数，则 $f$ 为满射。
\end{chapex}

\begin{chapex} % Always
设 $a,b,c,d \in \mathbb{R}$ 并定义函数 $f : \mathbb{R}^2 \to \mathbb{R}^2$ 为对于所有 $(x,y) \in \mathbb{R}^2$, $f(x,y) = (ax+by,cx+dy)$。则 $f$ 为单射当且仅当 $f$ 为满射。
\hintlabel{cqLinearMapInjectiveIffSurjective}{%
首先证明 $f$ 为双射当且仅当 $ad \ne bc$。
}
\end{chapex}

\begin{chapex} % Sometimes
\label{cqFunctionsASNEnd}
设 $U, V \subseteq \mathbb{R}$ 并假设对于所有 $x \in U$, $x^2 \in V$。函数 $f : U \to V$ 定义为 $f(x) = x^2$ 为单射。
\end{chapex}